% Source: Weinberg GR, physical PDF pp. 181--184
%         (printed pp. 157--160).
% Coverage: Section 7.3, The Brans-Dicke Theory; equations
%           (7.3.1)--(7.3.15).
% Figures/tables/footnotes: reference marker 2; no figures or tables.
% Status: source-reviewed and compile-clean.
% Uncertainties: none.
% MODERNIZATION: Curvature-dependent displays and their divergence derivation
% are converted to the target curvature convention and +8\pi/\phi source
% sign. Comma/semicolon notation is replaced by partial/nabla operators.

\subsection{The Brans-Dicke Theory}\label{sec:7.3}

Long-range forces are known to be transmitted by the gravitational field
\(g_{\mu\nu}\) and by the electromagnetic potential \(A_\mu\). It is natural
then to suspect that other long-range forces may be produced by scalar fields.
Such theories have been suggested since before general relativity; this
section describes the latest and possibly the best motivated theory in which
a scalar field shares the stage with gravitation, that of Brans and
Dicke.~\hyperref[ch7-ref-2]{[2]}

The starting point for Brans and Dicke is the idea of Mach, that the
phenomenon of inertia ought to arise from accelerations with respect to the
general mass distribution of the universe. (See Section~\ref{sec:1.3}.)
Thus the inertial masses of the various elementary particles ought not to be
fundamental constants, but should rather represent the particles' interaction
with some cosmic field. But the absolute scale of the elementary particle
masses (as opposed to their ratios, which presumably have nothing to do with
cosmic fields) can be measured only by measuring gravitational accelerations
\(Gm/r^2\), so an equivalent conclusion is that the gravitational constant
\(G\) ought to be related to the average value of a scalar field \(\phi\),
which is coupled to the mass density of the universe.

The simplest generally covariant field equation for such a scalar field would
be
\begin{equation}
  \Box\phi
  =
  4\pi\lambda T_M^\mu{}_\mu,
  \tag{7.3.1}\label{eq:7.3.1}
\end{equation}
where \(\Box\phi=\nabla_\mu\nabla^\mu\phi\) is now the invariant
d'Alembertian, \(\lambda\) is a coupling constant, and \(T_M^{\mu\nu}\) is
the energy-momentum tensor of the matter (i.e., everything but gravitation
and the \(\phi\) field) of the universe. We can make a rough estimate of the
average value of \(\phi\) by computing the central potential of a gas sphere
with the cosmic mass density \(\rho\sim10^{-29}\,\mathrm{g\,cm^{-3}}\) and
radius equal to the apparent radius of the universe
\(\mathcal R\sim10^{28}\,\mathrm{cm}\). (See Chapter~14.) This gives an
average value
\begin{equation}
  \langle\phi\rangle
  \sim
  \lambda\rho\mathcal R^2
  \sim
  \lambda\times10^{27}\,\mathrm{g\,cm^{-1}}.
  \tag{7.3.2}\label{eq:7.3.2}
\end{equation}
Note that \(10^{27}\,\mathrm{g\,cm^{-1}}\) is reasonably close to the
constant \(1/G=1.35\times10^{28}\,\mathrm{g\,cm^{-1}}\); hence we normalize
\(\phi\) so that
\begin{equation}
  \langle\phi\rangle\simeq\frac1G,
  \tag{7.3.3}\label{eq:7.3.3}
\end{equation}
and Eq.~\eqref{eq:7.3.2} then shows that \(\lambda\) is a dimensionless
number of order unity. These considerations led Brans and Dicke to suggest
that the correct field equations for gravitation are obtained by replacing
\(G\) with \(1/\phi\) and including an energy-momentum tensor
\(T_\phi^{\mu\nu}\) for the \(\phi\) field in the source of the
gravitational field:
\begin{equation}
  R^{\mu\nu}-\frac12g^{\mu\nu}R
  =
  \frac{8\pi}{\phi}
  \left(T_M^{\mu\nu}+T_\phi^{\mu\nu}\right).
  \tag{7.3.4}\label{eq:7.3.4}
\end{equation}

We do not, however, wish to give up the successes of the Principle of
Equivalence, such as the equality of gravitational and inertial mass, and the
gravitational time dilation. Brans and Dicke therefore require that it is
only \(g_{\mu\nu}\), and not \(\phi\), that enters in the equations of motion
of particles and photons. Therefore the equation describing the interchange
of energy between matter and gravitation is the same as in Einstein's theory:
\begin{equation}
  \nabla_\nu T_M^{\mu\nu}
  =
  \partial_\nu T_M^{\mu\nu}
  +\Gamma^\mu{}_{\nu\rho}T_M^{\rho\nu}
  +\Gamma^\nu{}_{\nu\rho}T_M^{\mu\rho}
  =
  0.
  \tag{7.3.5}\label{eq:7.3.5}
\end{equation}
The Bianchi identities tell us that the left-hand side of
Eq.~\eqref{eq:7.3.4} has vanishing covariant divergence, so by multiplying
Eq.~\eqref{eq:7.3.4} by \(\phi\) and taking the covariant divergence, we
find
\begin{equation}
  \left(R^\mu{}_\nu-\frac12\delta^\mu{}_\nu R\right)
  \nabla_\mu\phi
  =
  8\pi\nabla_\mu T_\phi^\mu{}_\nu.
  \tag{7.3.6}\label{eq:7.3.6}
\end{equation}

This requirement proves sufficient to determine \(T_\phi^\mu{}_\nu\). The
most general symmetric tensor that can be built up from terms each of which
involves two derivatives of one or two \(\phi\) fields, and \(\phi\) itself,
is
\begin{align}
  T_\phi^\mu{}_\nu
  ={}&
  A(\phi)\nabla^\mu\phi\nabla_\nu\phi
  +B(\phi)\delta^\mu{}_\nu
   \nabla_\rho\phi\nabla^\rho\phi
  \notag\\
  &+
  C(\phi)\nabla^\mu\nabla_\nu\phi
  +\delta^\mu{}_\nu D(\phi)\Box\phi.
  \tag{7.3.7}\label{eq:7.3.7}
\end{align}
A straightforward calculation gives
\begin{align}
  \nabla_\mu T_\phi^\mu{}_\nu
  ={}&
  [A'(\phi)+B'(\phi)]
  \nabla^\mu\phi\nabla_\nu\phi\nabla_\mu\phi
  \notag\\
  &+
  [A(\phi)+D'(\phi)]\nabla_\nu\phi\,\Box\phi
  \notag\\
  &+
  [A(\phi)+2B(\phi)+C'(\phi)]
  (\nabla^\mu\nabla_\nu\phi)\nabla_\mu\phi
  \notag\\
  &+
  D(\phi)\nabla_\nu\Box\phi
  +C(\phi)\Box(\nabla_\nu\phi).
  \tag{7.3.8}\label{eq:7.3.8}
\end{align}
(A prime here means the derivative with respect to \(\phi\).) The first term
of Eq.~\eqref{eq:7.3.6} is determined by Eq.~\eqref{eq:6.5.2} as
\begin{equation}
  \nabla_\sigma\phi\,R^\sigma{}_\nu
  =
  \Box(\nabla_\nu\phi)-\nabla_\nu\Box\phi.
  \tag{7.3.9}\label{eq:7.3.9}
\end{equation}
Also, by taking the trace of Eq.~\eqref{eq:7.3.4} and using
Eq.~\eqref{eq:7.3.1}, we find
\[
  R
  =
  -\frac{8\pi}{\phi}
  \left[
    \frac{1}{4\pi\lambda}\Box\phi
    +(A(\phi)+4B(\phi))\nabla_\mu\phi\nabla^\mu\phi
    +(C(\phi)+4D(\phi))\Box\phi
  \right],
\]
so the left-hand side of Eq.~\eqref{eq:7.3.6} is
\begin{align}
  \left(R^\mu{}_\nu-\frac12\delta^\mu{}_\nu R\right)
  \nabla_\mu\phi
  ={}&
  \Box(\nabla_\nu\phi)-\nabla_\nu\Box\phi
  \notag\\
  &+
  \frac{4\pi}{\phi}\nabla_\nu\phi
  \left[
    \left(
      \frac{1}{4\pi\lambda}
      +C(\phi)+4D(\phi)
    \right)\Box\phi
  \right.
  \notag\\[-2pt]
  &\hspace{6.2em}\left.
    +(A(\phi)+4B(\phi))
     \nabla_\mu\phi\nabla^\mu\phi
  \right].
  \tag{7.3.10}\label{eq:7.3.10}
\end{align}
By comparing the coefficients of
\(\nabla_\nu\Box\phi\), \(\Box(\nabla_\nu\phi)\),
\(\nabla_\nu\phi\,\Box\phi\),
\(\nabla^\mu\phi\nabla_\mu\phi\nabla_\nu\phi\), and
\((\nabla^\mu\nabla_\nu\phi)\nabla_\mu\phi\) in
Eqs.~\eqref{eq:7.3.8} and \eqref{eq:7.3.10}, we find that
Eq.~\eqref{eq:7.3.6} requires
\begin{align*}
  1&=-8\pi D(\phi),\\
  -1&=-8\pi C(\phi),\\
  -\frac{4\pi}{\phi}
  \left(\frac{1}{4\pi\lambda}+C(\phi)+4D(\phi)\right)
  &=-8\pi\left(A(\phi)+D'(\phi)\right),\\
  -\frac{4\pi}{\phi}\left(A(\phi)+4B(\phi)\right)
  &=-8\pi\left(A'(\phi)+B'(\phi)\right),\\
  0&=A(\phi)+2B(\phi)+C'(\phi).
\end{align*}
The unique solution is
\begin{equation}
  A(\phi)=\frac{\omega}{8\pi\phi},
  \qquad
  B(\phi)=-\frac{\omega}{16\pi\phi},
  \qquad
  C(\phi)=\frac1{8\pi},
  \qquad
  D(\phi)=-\frac1{8\pi},
  \tag{7.3.11}\label{eq:7.3.11}
\end{equation}
where \(\omega\) is a convenient dimensionless constant given by
\[
  \omega=\frac1\lambda-\frac32,
\]
or
\begin{equation}
  \lambda=\frac{2}{3+2\omega}.
  \tag{7.3.12}\label{eq:7.3.12}
\end{equation}
The field equations \eqref{eq:7.3.1} and \eqref{eq:7.3.4} of the
Brans--Dicke theory now read
\begin{equation}
  \Box\phi
  =
  \frac{8\pi}{3+2\omega}T_M^\mu{}_\mu,
  \tag{7.3.13}\label{eq:7.3.13}
\end{equation}
and
\begin{align}
  R_{\mu\nu}-\frac12g_{\mu\nu}R
  ={}&
  \frac{8\pi}{\phi}T_{M\mu\nu}
  +\frac{\omega}{\phi^2}
  \left(
    \nabla_\mu\phi\nabla_\nu\phi
    -\frac12g_{\mu\nu}\nabla_\rho\phi\nabla^\rho\phi
  \right)
  \notag\\
  &+
  \frac1\phi
  \left(
    \nabla_\mu\nabla_\nu\phi-g_{\mu\nu}\Box\phi
  \right).
  \tag{7.3.14}\label{eq:7.3.14}
\end{align}

Our previous estimate indicated that \(\lambda\) is of order unity, so we
expect that \(\omega\) is of order unity. If \(\omega\) is much larger than
unity, then Eq.~\eqref{eq:7.3.13} gives
\(\Box\phi=\order(1/\omega)\), and therefore
\begin{equation}
  \phi
  =
  \langle\phi\rangle+\order\left(\frac1\omega\right)
  =
  \frac1G+\order\left(\frac1\omega\right).
  \tag{7.3.15}\label{eq:7.3.15}
\end{equation}
Using this in Eq.~\eqref{eq:7.3.14} gives then
\[
  R_{\mu\nu}-\frac12g_{\mu\nu}R
  =
  8\pi G T_{M\mu\nu}
  +\order\left(\frac1\omega\right).
\]
Thus the Brans--Dicke theory goes over to the Einstein theory in the limit
\(\omega\to\infty\).

It must be stressed that the role of the scalar field in the Brans--Dicke
theory is confined to its effect on the gravitational field equations. Once
\(g_{\mu\nu}\) is calculated, the effects of gravitation on arbitrary
physical systems are to be determined exactly as described in Chapters~3
through~5.

Throughout most of this book it will be assumed that there is no scalar field
\(\phi\) that contributes to long-range interactions. However, from time to
time we return to the Brans--Dicke theory in order to see what changes it
would make in the predictions of general relativity.
